% CREATED FROM THESIS PROPOSAL AND IMPLEMENTATION NOTES
% Chapter 2: Theory and Background
% This file is intended to replace the template Theory.tex.
%
% Expected packages in the main thesis preamble:
% \usepackage{amsmath,amssymb}
% \usepackage{float}
% \usepackage{booktabs} % table style: top/mid/bottom rules, no vertical lines
% \usepackage{minted}

\chapter{Theory}
\label{ch:theory}

This chapter introduces the theoretical background for a runtime-level scheduling and DVFS-control framework for single-node multi-GPU task graphs. The target execution model is CUDA Sequential Task Flow (CUDASTF), where applications are represented as task graphs with data-driven dependencies. The central optimization goal is to reduce energy consumption and energy--delay product (EDP) while keeping the performance loss within a user-defined bound. To support this goal, the chapter first introduces multi-GPU task graph execution and CUDASTF, then reviews GPU dynamic voltage and frequency scaling (DVFS), energy--performance metrics, DAG scheduling, data-locality-aware placement, online bandit scheduling, and windowed proposal--commit DVFS control.

\section{Multi-GPU Task Graph Execution}
\label{sec:theory_task_graphs}

Task-graph execution is the foundation of this thesis. Instead of viewing a GPU application as a flat sequence of kernel launches, a task-based runtime represents the application as a graph of computation and data-movement tasks. This graph exposes dependencies, communication, device placement choices, and possible slack that can be exploited by a scheduler.

\subsection{Directed Acyclic Graph Representation}
\label{subsec:dag_representation}

A task graph is usually modelled as a directed acyclic graph (DAG):
\begin{equation}
G = (V,E),
\label{eq:dag_definition}
\end{equation}
where each vertex $v \in V$ represents a task and each directed edge $(u,v) \in E$ represents a dependency from task $u$ to task $v$. In a GPU runtime, a task may correspond to a kernel execution, a data transfer, a synchronization operation, or a higher-level runtime operation. A task $v$ can start only when all predecessor tasks in $\mathrm{pred}(v)$ have completed and the required input data are available on the device selected for $v$.

The acyclic property is important because it gives the runtime a partial order over tasks. Tasks that are not connected by dependencies may execute concurrently, while dependent tasks must respect the order encoded by the edges. This makes the DAG a natural abstraction for multi-GPU execution, where available parallelism and communication constraints must be considered together.

\subsection{Device Readiness, Data Readiness, and Finish Time}
\label{subsec:device_data_finish_time}

In a multi-GPU system, scheduling a task requires choosing both a device and a start time. Let $d$ denote a candidate GPU. The data-ready time of task $v$ on device $d$ can be expressed as
\begin{equation}
R_d(v) = \max_{u \in \mathrm{pred}(v)} \left( F(u) + C_{u,v,d} \right),
\label{eq:data_ready_time}
\end{equation}
where $F(u)$ is the finish time of predecessor $u$, and $C_{u,v,d}$ is the communication or data-migration cost needed before $v$ can execute on device $d$. If the required data are already resident on device $d$, then $C_{u,v,d}$ may be close to zero.

Let $A_d$ be the time at which device $d$ becomes available according to the scheduler. The estimated start time and finish time of task $v$ on device $d$ are then
\begin{align}
S_d(v) &= \max\left(A_d, R_d(v)\right),
\label{eq:start_time}\\
F_d(v) &= S_d(v) + T_d(v),
\label{eq:finish_time}
\end{align}
where $T_d(v)$ is the predicted execution time of task $v$ on device $d$. These equations are the basis for earliest-finish-time scheduling and for HEFT-style placement decisions.

\subsection{Critical Path, Slack, and Parallelism}
\label{subsec:critical_path_slack}

The critical path is the longest dependency chain in the DAG, including computation and communication costs. It determines the minimum possible makespan under ideal scheduling assumptions. Tasks on or near the critical path have little scheduling freedom because delaying them is likely to delay the whole application.

Slack is the opposite notion. A task has slack if it can be delayed, placed differently, or executed more slowly without extending the final makespan. Slack can arise from load imbalance, data transfers, synchronization, or non-critical branches of the DAG. In this thesis, slack is important for two reasons. First, non-critical tasks may tolerate energy-saving frequency choices. Second, communication and waiting periods may hide the latency of DVFS transitions.

\subsubsection{Why Slack Matters for DVFS}
\label{subsubsec:slack_for_dvfs}

DVFS reduces energy only when the saved power outweighs the additional runtime and transition overhead. If a task is on the critical path, lowering its frequency may directly increase the application runtime. If the same task has slack, the runtime increase may be partially or fully hidden. Therefore, a DAG-aware DVFS policy should not make frequency decisions from task cost alone; it should also consider dependency structure, data movement, and criticality.

\section{CUDASTF Execution Model}
\label{sec:theory_cudastf}

CUDASTF is the runtime substrate considered in this thesis. It implements a CUDA-oriented Sequential Task Flow model in which applications are expressed through task submissions with data-access information. The runtime derives data-driven dependencies, manages data placement, and schedules work across GPUs \cite{Augonnet2024CUDASTF}.

\subsection{Sequential Task Flow Model}
\label{subsec:stf_model}

The Sequential Task Flow model allows the programmer to submit tasks sequentially while the runtime extracts parallelism from data dependencies. The key idea is that the programmer specifies which data objects a task reads or writes. The runtime then determines when the task can execute safely. If two tasks access independent data, they may run concurrently. If one task writes data that another task later reads, the runtime inserts the required dependency.

This model is useful for GPU applications because explicit synchronization is often error-prone and overly conservative. A task-based runtime can avoid unnecessary synchronization while preserving correctness through data-dependency tracking.

\subsection{Data-Driven Dependencies and Data Residency}
\label{subsec:cudastf_data_dependencies}

In multi-GPU execution, dependencies are not only ordering constraints. They also determine where data are located and whether a task requires data migration. CUDASTF tracks data accesses and can therefore infer when a task needs a local copy, when data must be transferred between GPUs, and when multiple devices may share valid copies.

For the scheduler, this information is essential. A device that appears fast from a compute perspective may be a poor choice if the required data are located elsewhere. Conversely, a slightly slower or more loaded device may be attractive if it already holds the necessary input data.

\subsection{Runtime Scheduling Opportunities}
\label{subsec:cudastf_scheduling_opportunities}

CUDASTF exposes several opportunities for energy-aware scheduling:
\begin{itemize}
    \item task dependencies can be used to estimate criticality and slack;
    \item data residency information can be used to reduce unnecessary transfers;
    \item communication phases can be used to hide frequency-transition latency;
    \item repeated task types can be profiled and modelled at runtime;
    \item task placement and DVFS decisions can be integrated into one runtime-level control pipeline.
\end{itemize}

These properties make CUDASTF a suitable setting for combining DAG scheduling with GPU DVFS. The thesis does not treat DVFS as an external node-level controller. Instead, DVFS decisions are made with direct access to task-graph information.

\section{GPU DVFS and Energy--Performance Trade-offs}
\label{sec:theory_dvfs}

Dynamic voltage and frequency scaling (DVFS) is a hardware power-management mechanism that changes operating frequencies, and when supported, the corresponding voltages. On GPUs, DVFS is commonly exposed through core and memory frequency controls. Prior GPU DVFS studies show that the best frequency setting depends strongly on workload characteristics, and that maximum frequency is not always energy-optimal \cite{Mei2017Survey,Ali2022Optimal,Ali2023Portable,Wang2024DSO}.

\subsection{GPU Core and Memory Frequency Domains}
\label{subsec:core_memory_frequency}

Two frequency domains are particularly relevant for this thesis. The GPU core frequency affects arithmetic throughput, instruction issue, and the execution rate of compute pipelines. The memory frequency affects memory bandwidth and, indirectly, the cost of moving data to and from global memory.

A frequency configuration can be written as
\begin{equation}
\phi = (f_g, f_m),
\label{eq:frequency_config}
\end{equation}
where $f_g$ is the GPU core frequency and $f_m$ is the memory frequency. A DVFS policy chooses $\phi$ for a task, a group of tasks, or a larger execution phase.

Although the hardware may expose many supported clock pairs, the runtime policy in this thesis uses a small discrete abstraction:
\begin{equation}
L = \{\mathrm{HIGH},\mathrm{MID},\mathrm{LOW}\}.
\label{eq:dvfs_levels}
\end{equation}
The high level represents the fastest available graphics clock, while the mid and low levels are selected by rank among lower supported graphics clocks. The memory clock is treated conservatively and kept at the highest supported level in the thesis design. This discrete-level view keeps the online decision space small enough for a runtime scheduler while still exposing meaningful energy--performance choices.

\subsection{Workload Sensitivity to Frequency Scaling}
\label{subsec:workload_sensitivity}

The effect of DVFS depends on how a task uses the GPU. The thesis uses the following conceptual classification.

\subsubsection{Compute-Bound Tasks}
\label{subsubsec:compute_bound}

A compute-bound task spends most of its time in arithmetic operations or instruction execution. Its runtime is sensitive to GPU core frequency. Lowering the core frequency may reduce power, but it is also likely to increase runtime. For compute-bound tasks on the critical path, aggressive down-scaling can therefore harm EDP.

\subsubsection{Memory-Bound Tasks}
\label{subsubsec:memory_bound}

A memory-bound task spends a large fraction of time waiting for memory accesses. Such a task may be less sensitive to core frequency because increasing arithmetic throughput does not remove the memory bottleneck. In these cases, lower core frequency may save energy with limited performance impact. However, memory frequency may still be important.

\subsubsection{Mixed Tasks}
\label{subsubsec:mixed_tasks}

Many GPU tasks are neither purely compute-bound nor purely memory-bound. Their sensitivity depends on input size, data layout, reuse, occupancy, and concurrent runtime activities. Mixed tasks motivate lightweight runtime models rather than fixed, hand-written rules for all kernels.

\subsection{Energy of a Frequency Configuration}
\label{subsec:energy_frequency_config}

For a task $v$ executed under frequency configuration $\phi$, the task energy can be approximated as
\begin{equation}
E_v(\phi) = P_v(\phi) \cdot T_v(\phi),
\label{eq:task_energy}
\end{equation}
where $P_v(\phi)$ is average power and $T_v(\phi)$ is execution time. The energy-optimal configuration is
\begin{equation}
\phi_v^{*} = \arg\min_{\phi \in \Phi} E_v(\phi),
\label{eq:energy_optimal_frequency}
\end{equation}
where $\Phi$ is the set of supported frequency configurations.

However, minimizing per-task energy independently is not necessarily optimal at the application level. A task-level decision may increase the critical path, create queueing effects, or trigger excessive frequency transitions. Therefore, the thesis focuses on runtime-level decisions that consider both task characteristics and DAG structure.

\subsection{DVFS Transition Latency}
\label{subsec:dvfs_transition_latency}

Changing GPU frequency is not free. A transition may introduce latency, temporarily stall command execution, or interact with stream scheduling. If the runtime changes frequency before every short task, the accumulated transition cost can dominate the potential savings. A more practical approach is to amortize transitions over groups of tasks and to place transitions where their latency is least visible, such as during data transfers or non-critical work.

\section{Energy, EDP, and Bounded Slowdown Objectives}
\label{sec:theory_metrics}

The scheduler in this thesis is multi-objective. It should reduce energy consumption while preserving acceptable performance. This section defines the metrics used to reason about this trade-off.

\subsection{Runtime, Power, and Energy}
\label{subsec:runtime_power_energy}

The application runtime is the makespan of the task graph:
\begin{equation}
T_{\mathrm{app}} = \max_{v \in V} F(v).
\label{eq:application_runtime}
\end{equation}

The total energy is the integral of power over time:
\begin{equation}
E_{\mathrm{app}} = \int_0^{T_{\mathrm{app}}} P(t)\,dt.
\label{eq:application_energy_integral}
\end{equation}
In an experimental system, this integral is usually approximated from sampled power readings:
\begin{equation}
E_{\mathrm{app}} \approx \sum_i P_i \Delta t_i.
\label{eq:application_energy_sampled}
\end{equation}

\subsection{Energy--Delay Product and Related Metrics}
\label{subsec:edp_metrics}

Energy--delay product combines energy and runtime:
\begin{equation}
\mathrm{EDP} = E_{\mathrm{app}} \cdot T_{\mathrm{app}}.
\label{eq:edp}
\end{equation}
A more performance-sensitive form is
\begin{equation}
\mathrm{ED}^{2}\mathrm{P} = E_{\mathrm{app}} \cdot T_{\mathrm{app}}^2.
\label{eq:ed2p}
\end{equation}

EDP is useful when both energy and runtime matter. ED$^2$P penalizes runtime degradation more strongly and is therefore useful when performance loss must be kept small.

\subsection{Bounded Slowdown}
\label{subsec:bounded_slowdown}

A bounded slowdown constraint compares the optimized runtime to a baseline runtime $T_{\mathrm{base}}$:
\begin{equation}
T_{\mathrm{app}} \leq (1 + \epsilon) T_{\mathrm{base}},
\label{eq:bounded_slowdown}
\end{equation}
where $\epsilon$ is the maximum allowed slowdown. The baseline may be fixed maximum-frequency execution or another user-selected reference policy.

\begin{table}[H]
\centering
\begin{tabular}{@{}ll@{}}
\toprule
Metric & Definition \\
\midrule
Runtime & $T_{\mathrm{app}}$ \\
Energy & $E_{\mathrm{app}}$ \\
EDP & $E_{\mathrm{app}}T_{\mathrm{app}}$ \\
ED$^2$P & $E_{\mathrm{app}}T_{\mathrm{app}}^2$ \\
Slowdown & $T_{\mathrm{app}}/T_{\mathrm{base}} - 1$ \\
\bottomrule
\end{tabular}
\caption{Energy--performance metrics used in this thesis.}
\label{tab:energy_performance_metrics}
\end{table}

\section{Online Task Characterization for DVFS}
\label{sec:theory_task_characterization}

A runtime scheduler cannot rely on expensive offline modelling for every task instance. Offline task-type profiles can be useful background information, but the DVFS path studied in this thesis is designed to operate from information already available inside the scheduler. The characterization problem is therefore not to predict exact power and runtime for every frequency pair. It is to separate tasks that should remain at high frequency from tasks that are plausible mid- or low-frequency candidates under the current DAG and device context.

\subsection{Observable Runtime Features}
\label{subsec:observable_runtime_features}

For each scheduled task $v$, the runtime observes or estimates a small set of features:
\begin{itemize}
    \item task cost $T(v)$;
    \item dependency fan-in and input-byte volume;
    \item copy time and queue-wait time under the selected placement context;
    \item device backlog on the chosen GPU;
    \item a local criticality proxy derived from cost, dependencies, and input size;
    \item task-symbol hints for heavy kernels such as dense linear algebra operations.
\end{itemize}
These features are cheap to compute because they are already needed for HEFT-style candidate evaluation or are available from CUDASTF task metadata.

\subsection{Copy-Wait Ratio and Slack Proxy}
\label{subsec:copy_wait_slack_proxy}

The DVFS model uses copy and waiting time as a proxy for how much execution can hide slower core frequency. Let $W(v)$ denote the predicted copy-plus-queue waiting component and let $T(v)$ denote the predicted compute component. A normalized copy-wait ratio is
\begin{equation}
\omega(v) =
\frac{W(v)}{\max(\varepsilon, W(v)+T(v))}.
\label{eq:copy_wait_ratio}
\end{equation}
A large $\omega(v)$ indicates that the task is dominated by data movement or queue waiting, so a lower core frequency is less likely to be fully visible in the application makespan. The scheduler also uses a slack proxy
\begin{equation}
s(v)=1-\kappa(v),
\label{eq:slack_proxy}
\end{equation}
where $\kappa(v)$ is the local criticality proxy. This is not an exact critical-path slack computation, but it gives the DVFS layer a conservative signal: high-criticality tasks are treated as having little slack, while low-criticality tasks may be considered for lower frequency when other guard conditions agree.

\subsection{Runtime Suitability}
\label{subsec:runtime_suitability}

The model must be inexpensive to evaluate and robust to imperfect estimates. For this reason, the thesis uses online proxies and guardrails rather than a large offline table of per-task frequency sensitivities. A task may request lower frequency only when multiple signals agree: it is not critical, not urgently heavy, not on a backlogged device, and has enough copy-wait time or slack to hide the expected slowdown. This conservative characterization supports the design goal of exposing an energy-control path without turning every task into a separate profiling problem.

\section{HEFT-Based DAG Scheduling}
\label{sec:theory_heft}

Heterogeneous Earliest Finish Time (HEFT) is a widely used list-scheduling heuristic for heterogeneous systems \cite{Topcuoglu2002HEFT}. The full HEFT algorithm ranks tasks and then places them on processors according to earliest finish time. This thesis uses HEFT primarily as a baseline placement principle: for each ready task, estimate its finish time on each GPU and select the device with the earliest finish time.

\subsection{Earliest-Finish-Time Placement}
\label{subsec:eft_placement}

Given the candidate finish time $F_d(v)$ from Equation~\eqref{eq:finish_time}, a HEFT-style baseline chooses
\begin{equation}
 d_{\mathrm{base}}(v) = \arg\min_{d \in D} F_d(v),
\label{eq:heft_baseline_device}
\end{equation}
where $D$ is the set of available GPUs. This decision balances device availability, data-readiness, communication cost, and predicted execution time.

\subsection{Strengths of HEFT for Multi-GPU Runtime Scheduling}
\label{subsec:heft_strengths}

HEFT-style scheduling is attractive because it is simple, fast, and directly connected to makespan minimization. It naturally avoids placing tasks on devices that are busy or require excessive data movement if doing so would delay the finish time. For online runtime scheduling, this makes HEFT a strong baseline.

\subsection{Limitations for Energy-Aware Scheduling}
\label{subsec:heft_limitations}

The main limitation is that HEFT is performance-oriented. It does not directly optimize energy, EDP, or DVFS transition overhead. It may also ignore placement alternatives that slightly delay a non-critical task but save a significant amount of data movement. This motivates a controlled-deviation strategy: keep HEFT as the safe baseline, but allow limited deviations when the expected communication benefit is high and the end-time penalty is small.

\section{Data Locality, Communication Cost, and Criticality}
\label{sec:theory_data_locality}

Multi-GPU applications are often limited not only by computation but also by data movement. A scheduler that ignores data locality may repeatedly move data between GPUs, increasing runtime and energy. This section explains the concepts used to reason about locality-aware placement.

\subsection{Data Residency and Copy Cost}
\label{subsec:data_residency_copy_cost}

Data residency describes which GPU currently holds a valid copy of a data object. If a task is placed on a GPU that already contains its required inputs, it can start after dependency completion and device availability. If the data are located elsewhere, the runtime must perform one or more transfers.

For a candidate device $d$, the total copy cost for task $v$ can be approximated as
\begin{equation}
C_d(v) = \sum_{x \in \mathrm{inputs}(v)} C_{x,d},
\label{eq:copy_cost}
\end{equation}
where $C_{x,d}$ is the estimated cost of making input data object $x$ available on device $d$.

\subsection{Queue Delay and Data-Ready Delay}
\label{subsec:queue_data_delay}

A placement candidate can be delayed for two different reasons. First, the selected device may already be busy. Second, the required data may not yet be available. These two delays can be separated as
\begin{align}
Q_d(v) &= \max\left(0, S_d(v) - R_d(v)\right), \\
W_d(v) &= \max\left(0, R_d(v) - A_d\right),
\end{align}
where $Q_d(v)$ represents queue delay after data are ready, and $W_d(v)$ represents waiting for data relative to device availability. This distinction helps identify whether a task is delayed primarily by load imbalance or by data movement.

\subsection{Criticality Proxy}
\label{subsec:criticality_proxy}

Exact critical-path computation can be expensive or unavailable in online runtime scheduling. A lightweight proxy can estimate whether a task is likely to be important. The current scheduling design uses three normalized features: task cost, dependency count, and input size. Let
\begin{equation}
c_T(v),\qquad c_D(v),\qquad c_B(v) \in [0,1]
\end{equation}
denote the normalized task-cost, dependency-count, and input-byte features for task $v$. The criticality proxy combines them with tunable non-negative weights:
\begin{equation}
\kappa(v) = \mathrm{clip}_{[0,1]}\left(
\beta_C c_T(v) + \beta_d c_D(v) + \beta_B c_B(v)
\right),
\label{eq:criticality_proxy}
\end{equation}
where
\begin{equation}
\beta_C,\beta_d,\beta_B \geq 0,
\qquad
\beta_C+\beta_d+\beta_B=1.
\end{equation}
The weights describe how much the proxy should trust each signal. A larger $\beta_C$ makes long-running tasks more critical; a larger $\beta_d$ emphasizes tasks with many dependencies, which often occur near synchronization points; and a larger $\beta_B$ emphasizes data-heavy tasks, where poor placement can amplify communication cost. In the parameterization used in this thesis, the default values are
\begin{equation}
\beta_C=0.55,\qquad \beta_d=0.30,\qquad \beta_B=0.15.
\end{equation}
These constants are therefore not part of the theoretical definition itself. They are a default parameterization of a weighted proxy whose principle is that long, dependency-heavy, data-heavy tasks should be treated more conservatively.

\section{Controlled Deviation from HEFT Baselines}
\label{sec:theory_controlled_deviation}

This thesis does not discard the HEFT baseline. Instead, it uses HEFT as a conservative placement decision and then considers whether a nearby alternative is worthwhile. The goal is to reduce communication and improve locality without significantly extending the task graph's critical path.

\subsection{Baseline and Alternative Candidates}
\label{subsec:baseline_alt_candidates}

For each task $v$, the scheduler first computes the baseline device $d_{\mathrm{base}}$ using Equation~\eqref{eq:heft_baseline_device}. It then evaluates alternative devices $d \neq d_{\mathrm{base}}$. An alternative is interesting when it saves copy cost or queue delay while keeping the predicted finish time close to the baseline.

\subsection{Copy Gain and End-Time Penalty}
\label{subsec:copy_gain_end_penalty}

Let $F_{\mathrm{base}}(v)$ and $C_{\mathrm{base}}(v)$ be the baseline finish time and copy cost. For an alternative device $d$, define the relative end-time penalty as
\begin{equation}
\Delta_F(d,v) = \frac{F_d(v) - F_{\mathrm{base}}(v)}{\max(\varepsilon, F_{\mathrm{base}}(v))},
\label{eq:end_penalty}
\end{equation}
where $\varepsilon$ avoids division by zero. The copy gain is
\begin{equation}
G_C(d,v) = \frac{C_{\mathrm{base}}(v) - C_d(v)}{\max(C_{\min}, C_{\mathrm{base}}(v))},
\label{eq:copy_gain}
\end{equation}
where $C_{\min}$ is a small denominator floor.

\subsection{Gate-Based Acceptance}
\label{subsec:gate_based_acceptance}

An alternative is accepted only if it satisfies a gate:
\begin{align}
\Delta_F(d,v) &\leq \tau_F, \\
G_C(d,v) &\geq \tau_C,
\label{eq:residual_gate}
\end{align}
where $\tau_F$ is the allowed end-time penalty and $\tau_C$ is the required copy-gain threshold. Critical tasks should use stricter thresholds than non-critical tasks.

\subsubsection{Critical and Non-Critical Gates}
\label{subsubsec:critical_noncritical_gates}

The scheduler can use two sets of thresholds. For critical tasks, the allowed end-time penalty should be small and the copy-gain requirement should be high. For non-critical tasks, the scheduler can be more aggressive because small local delays may be hidden by slack.

\subsubsection{Gate Profiles}
\label{subsubsec:gate_profiles}

A gate profile is a parameterized policy that controls how willing the scheduler is to deviate from HEFT. Table~\ref{tab:gate_profiles_theory} gives a conceptual example.

\begin{table}[H]
\centering
\begin{tabular}{@{}llll@{}}
\toprule
Profile & Behaviour & End-time penalty & Copy-gain threshold \\
\midrule
Closed & Always follow HEFT & Not applicable & Not applicable \\
Conservative & Rarely deviate & Very small & High \\
Default & Balanced deviation & Small & Moderate \\
Aggressive & More locality-driven & Larger & Lower \\
\bottomrule
\end{tabular}
\caption{Conceptual gate profiles for controlled deviation from HEFT.}
\label{tab:gate_profiles_theory}
\end{table}

This profile abstraction is useful because an online scheduler can adapt the profile according to recent workload behaviour instead of relying on one fixed policy.

\section{Contextual Bandits for Online Runtime Scheduling}
\label{sec:theory_contextual_bandits}

Runtime scheduling is online: the scheduler must make decisions with incomplete information and adapt as execution progresses. Contextual bandits provide a lightweight framework for choosing among several policies while learning from observed rewards. In this thesis, the bandit does not select a GPU directly. Instead, it selects a gate profile that controls how aggressively the scheduler may deviate from the HEFT baseline.

\subsection{Exploration and Exploitation}
\label{subsec:exploration_exploitation}

A scheduler faces an exploration--exploitation trade-off. Exploitation means using the currently best-known profile. Exploration means trying another profile to learn whether it performs better under the current workload. Pure exploitation may get stuck in a suboptimal policy, while excessive exploration may harm performance.

\subsection{Window-Level Context}
\label{subsec:window_level_context}

Instead of choosing a new bandit action for every task, the scheduler can operate at window granularity. A window contains a fixed number of tasks. At the beginning of a window, the scheduler observes a context vector $x$ summarizing recent execution, such as average criticality, average copy cost, average task cost, device load imbalance, and recent deviation rate. It then selects a gate profile for all tasks in that window.

This window-level design reduces noise and overhead. It also makes reward attribution easier because the reward is aggregated over many task decisions.

\subsection{LinUCB Selection}
\label{subsec:linucb_selection}

LinUCB is a contextual bandit algorithm that assumes the expected reward of an action is approximately linear in the context \cite{Li2010ContextualBandit}. For each arm $a$, the algorithm maintains a parameter estimate $\theta_a$. Given context $x$, it computes
\begin{equation}
\mathrm{score}_a = \theta_a^{T}x + \alpha \sqrt{x^{T}A_a^{-1}x},
\label{eq:linucb_score}
\end{equation}
where the first term is the predicted reward, the second term is an uncertainty bonus, $A_a^{-1}$ is the inverse covariance matrix for arm $a$, and $\alpha$ controls exploration. The scheduler selects the arm with the highest score.

\subsection{Reward for Placement Profiles}
\label{subsec:placement_reward}

A useful reward should encourage copy savings while penalizing risky delays. For a deviated task, a simple reward can be written as
\begin{equation}
r(v) = G_C(v) - \lambda_F \Delta_F(v),
\label{eq:task_reward}
\end{equation}
where $G_C(v)$ is copy gain, $\Delta_F(v)$ is end-time penalty, and $\lambda_F$ is larger for critical tasks than for non-critical tasks.

At the end of a window $w$, the reward can be averaged and penalized by the deviation rate:
\begin{equation}
R_w = \frac{1}{|w|}\sum_{v \in w} r(v) - \eta \rho_w,
\label{eq:window_reward}
\end{equation}
where $\rho_w$ is the fraction of tasks that deviated from HEFT and $\eta$ is a regularization coefficient. This avoids rewarding a policy that deviates too often without sufficient benefit.

\subsubsection{Relative Reward}
\label{subsubsec:relative_reward}

A relative reward compares the selected profile against a default profile on the same window or task. This reduces the effect of workload difficulty. For example, if no profile has good deviation opportunities in a window, the reward should not incorrectly punish the selected profile. Relative reward therefore helps the bandit learn whether a profile is better than the default under the same scheduling conditions.

\section{Region-Level DVFS and Transition Amortization}
\label{sec:theory_region_dvfs}

DVFS can be applied at different granularities. Very fine-grained control may choose a frequency for every kernel, while coarse-grained control may choose one frequency for a whole application phase. This thesis uses the task graph to define an intermediate granularity: frequency regions or device-level commit windows.

\subsection{DVFS Granularity}
\label{subsec:dvfs_granularity}

Table~\ref{tab:dvfs_granularity} summarizes common DVFS granularities.

\begin{table}[H]
\centering
\begin{tabular}{@{}lll@{}}
\toprule
Granularity & Advantage & Limitation \\
\midrule
Per-kernel / per-task & Fine control & Many transitions, high overhead \\
Region / window & Amortizes transitions & Requires grouping and aggregation \\
Application phase & Low overhead & May miss fine-grained variation \\
Whole application & Simple and stable & Too coarse for heterogeneous DAGs \\
\bottomrule
\end{tabular}
\caption{DVFS granularity levels and their trade-offs.}
\label{tab:dvfs_granularity}
\end{table}

Per-task DVFS can approximate the best frequency for each task but may trigger too many transitions. Phase-level DVFS is stable but may ignore short-term changes in task type, data movement, and GPU load. Region-level DVFS is a compromise: it allows the runtime to react to task-graph structure while limiting transition overhead.

\subsection{Frequency Regions in Task Graphs}
\label{subsec:frequency_regions}

A frequency region is a group of tasks that share a frequency decision. Regions may be constructed along dependency chains, within a device window, or around repeated task types. The goal is to place tasks with similar frequency preferences into the same region and to avoid unnecessary frequency switches.

For a region $R$, a frequency decision can be written as
\begin{equation}
\phi_R^{*} = \arg\min_{\phi \in \Phi} J_R(\phi),
\label{eq:region_frequency_decision}
\end{equation}
where $J_R$ is an objective such as energy, EDP, or energy under a bounded slowdown constraint. In practice, the runtime may not solve this optimization exactly; it can use heuristic rules or learned policies that approximate the same goal.

\subsection{Overlapping Transitions with Communication and Slack}
\label{subsec:overlap_transition_slack}

Let $L_{\mathrm{tr}}$ be the latency of a frequency transition. If the transition is placed on the critical path, it may increase the makespan by nearly $L_{\mathrm{tr}}$. If it is placed during a data transfer or a period of slack, the visible delay can be smaller:
\begin{equation}
L_{\mathrm{visible}} = \max(0, L_{\mathrm{tr}} - L_{\mathrm{hidden}}),
\label{eq:visible_transition_latency}
\end{equation}
where $L_{\mathrm{hidden}}$ is the portion overlapped with communication, synchronization, or non-critical work.

This is one of the main reasons to integrate DVFS with task-graph scheduling. A runtime that sees dependencies and data movement can choose transition points more intelligently than an external controller that only observes coarse application phases.

\section{Windowed Proposal--Commit DVFS Control}
\label{sec:theory_proposal_commit}

The DVFS design in this thesis separates task-level frequency intent from device-level frequency commitment. This avoids making the placement bandit responsible for a very large joint action space of devices and frequencies. Instead, placement is the first layer, and DVFS is the second layer, with committed frequency selected at a per-device window granularity rather than independently for every task.

\subsection{Two-Layer Runtime Control}
\label{subsec:two_layer_control}

The runtime control problem is separated into two decisions:
\begin{enumerate}
    \item \textbf{Placement layer}: choose the GPU for each task, using a HEFT baseline and controlled deviation.
    \item \textbf{DVFS layer}: choose the frequency level for the selected GPU, using proposal--commit control.
\end{enumerate}

This separation simplifies learning and attribution. The placement layer reasons about data locality, queue delay, and finish time. The DVFS layer reasons about criticality, slack, backlog, and frequency risk after the device has been selected. It also makes the decisions easier to audit: placement decisions can be explained relative to HEFT, while frequency decisions can be explained through proposal shares and commit guardrails.

\subsection{Task-Level Proposer}
\label{subsec:task_level_proposer}

The proposer assigns each scheduled task a frequency intent, such as high, mid, or low. The intent is not immediately applied to that same task. It is a vote that expresses what the task appears to need under the current context and contributes to the next device-level commit decision. The current task receives the frequency level that was already committed by the previous window for the selected GPU.

A proposer can use the following features:
\begin{itemize}
    \item criticality proxy $\kappa(v)$;
    \item estimated task cost $T(v)$;
    \item backlog pressure on the selected device;
    \item copy/wait ratio;
    \item available slack;
    \item whether the task is heavy or latency-sensitive.
\end{itemize}

A conceptual scoring function is
\begin{equation}
s(v) = w_{\kappa}\kappa(v) + w_T c_T(v) + w_B b(v)
       - w_W \omega(v) - w_S s_{\mathrm{slack}}(v) + w_H h(v),
\label{eq:dvfs_proposer_score}
\end{equation}
where $b(v)$ is backlog pressure, $\omega(v)$ is the copy-wait ratio from Equation~\eqref{eq:copy_wait_ratio}, $s_{\mathrm{slack}}(v)$ is normalized slack, and $h(v)$ indicates a heavy task. Positive terms push the decision toward high frequency, while negative terms indicate that lower frequency may be safe.

\subsubsection{High-Frequency Conditions}
\label{subsubsec:high_frequency_conditions}

A task should request high frequency if it is critical, heavy with little slack, placed on a backlogged device, or expected to be highly sensitive to frequency reduction. High frequency is also the safe fallback when the runtime lacks enough information.

\subsubsection{Mid- and Low-Frequency Conditions}
\label{subsubsec:mid_low_conditions}

A task may request mid frequency when it has sufficient copy/wait time, moderate or low criticality, and no urgent heavy-computation requirement. Low frequency should be more restrictive. It is appropriate only for light, non-critical tasks with substantial slack and low backlog pressure.

\subsection{Device-Level Window Committer}
\label{subsec:device_window_committer}

The committer aggregates task proposals over a device-level window. Only the committer changes the actual committed frequency level, and the new level affects later tasks on the same GPU. This delayed-commit design prevents individual tasks from causing frequent frequency oscillations and amortizes transition overhead across a region of execution.

For a device window $W_d$, weighted proposal shares can be written as
\begin{align}
P_{\mathrm{high}} &= \frac{\sum_{v \in W_d} \omega_v \mathbb{I}[p_v=\mathrm{high}]}{\sum_{v \in W_d} \omega_v}, \\
P_{\mathrm{mid}} &= \frac{\sum_{v \in W_d} \omega_v \mathbb{I}[p_v=\mathrm{mid}]}{\sum_{v \in W_d} \omega_v}, \\
P_{\mathrm{low}} &= \frac{\sum_{v \in W_d} \omega_v \mathbb{I}[p_v=\mathrm{low}]}{\sum_{v \in W_d} \omega_v},
\label{eq:proposal_shares}
\end{align}
where $p_v$ is the proposal of task $v$ and $\omega_v$ is a proposal weight. Critical and expensive tasks can receive larger weights.

\subsection{Commit Eligibility and Slack Budget}
\label{subsec:commit_eligibility}

The committer should enter a lower frequency only when the estimated performance risk fits within the available slack. A simple estimate of extra runtime when moving from high frequency to target frequency $\phi$ is
\begin{equation}
\Delta T_v(\phi) = \alpha_v T_v \left(1 - \omega(v)\right)
\left(\frac{f_{g,\mathrm{high}}}{f_{g,\phi}} - 1\right),
\label{eq:estimated_extra_runtime}
\end{equation}
where $\omega(v)$ is the copy-wait ratio and $\alpha_v$ increases for critical or heavy tasks. The aggregate extra runtime should satisfy a slack-budget condition:
\begin{equation}
\sum_{v \in W_d} \Delta T_v(\phi) \leq \beta S_{W_d} + B,
\label{eq:slack_budget_condition}
\end{equation}
where $S_{W_d}$ is the estimated slack reserve, $\beta$ controls how much slack can be consumed, and $B$ is a small performance budget.

\subsection{Guardrails and Stability}
\label{subsec:guardrails_stability}

A proposal--commit system needs guardrails to avoid performance regressions. Table~\ref{tab:dvfs_guardrails} summarizes the guardrails used by the windowed controller.
\begin{table}[t]
\centering
\begin{tabular}{@{}p{0.26\linewidth}p{0.34\linewidth}p{0.32\linewidth}@{}}
\toprule
Guardrail & Trigger & Effect \\
\midrule
Critical-work protection & High criticality mass or urgent heavy tasks in the window & Force the next committed level to $\mathrm{HIGH}$ \\
Backlog protection & High backlog pressure on the selected GPU & Prevent lower-frequency commits while the device is overloaded \\
Frequency-separation check & Mid or low graphics clocks are too close to the high clock & Treat the lower level as unavailable to avoid ineffective transitions \\
Eligibility persistence & Mid or low eligibility appears only transiently & Require consecutive eligible windows before entering a lower level \\
Slowdown recovery & Recent window runtime or busy-tail time exceeds moving-average guard bands & Force recovery windows at $\mathrm{HIGH}$ \\
Low-level cooldown & A slowdown is observed after a low-frequency commit & Temporarily disable $\mathrm{LOW}$ while keeping $\mathrm{MID}$ available if safe \\
\bottomrule
\end{tabular}
\caption{Guardrails used by the windowed proposal--commit DVFS controller.}
\label{tab:dvfs_guardrails}
\end{table}

The purpose of these guardrails is not to maximize energy savings for every local window. Instead, they protect the global runtime objective and make DVFS behaviour stable. Recent slowdown is detected from window-level runtime and busy-tail measurements maintained as moving averages. If enough recent windows exceed the allowed slowdown ratios, the committer forces recovery to high frequency. This turns the lower-frequency levels into optimistic choices that must continue to earn their place.

\subsection{Algorithm}
\label{subsec:algorithm_proposal_commit}

Algorithm~\ref{alg:proposal_commit_pipeline} formalizes one scheduling step with windowed proposal--commit DVFS.

\begin{algorithm}[t]
\caption{Runtime scheduling with windowed proposal--commit DVFS}
\label{alg:proposal_commit_pipeline}
\begin{algorithmic}[1]
\Require Ready task $v$, GPU set $D$, placement state $\mathcal{S}$, DVFS windows $\{W_d\}_{d\in D}$
\Ensure Selected GPU $d_{\mathrm{exec}}$ and frequency annotation for task $v$
\ForAll{$d \in D$}
    \State $\xi(v,d) \gets \Call{EvaluateCandidate}{v,d,\mathcal{S}}$
\EndFor
\State $d_{\mathrm{base}} \gets \arg\min_{d\in D} F_d(v)$
\State $d_{\mathrm{exec}} \gets \Call{ResidualGate}{v,d_{\mathrm{base}},\{\xi(v,d)\}_{d\in D}}$
\State $\kappa(v) \gets \Call{CriticalityProxy}{v}$
\State $p_v \gets \Call{ProposeLevel}{v,d_{\mathrm{exec}},\kappa(v),\xi(v,d_{\mathrm{base}})}$
\State $\Call{AccumulateIntent}{W_{d_{\mathrm{exec}}},p_v,v}$
\State $\ell_v \gets \Call{CommittedLevel}{W_{d_{\mathrm{exec}}}}$
\State $\varphi_v \gets \Call{FrequencyPair}{d_{\mathrm{exec}},\ell_v}$
\State $\Call{AnnotateFrequency}{v,\varphi_v}$
\State $\mathcal{S} \gets \Call{UpdatePlacementState}{\mathcal{S},v,d_{\mathrm{exec}}}$
\If{$\Call{WindowComplete}{W_{d_{\mathrm{exec}}}}$}
    \State $m \gets \Call{AggregateWindowMetrics}{W_{d_{\mathrm{exec}}}}$
    \State $\ell_{\mathrm{next}} \gets \Call{CommitLevelWithGuardrails}{m}$
    \State $\Call{ResetWindow}{W_{d_{\mathrm{exec}}},\ell_{\mathrm{next}}}$
\EndIf
\State \Return $(d_{\mathrm{exec}},\ell_v)$
\end{algorithmic}
\end{algorithm}

\section{Multi-Objective Runtime Scheduling Framework}
\label{sec:theory_runtime_framework}

The preceding sections describe the components required by the thesis framework. This section summarizes how they fit together.

\subsection{Joint Placement and Frequency Control}
\label{subsec:joint_placement_frequency}

The runtime simultaneously considers device placement and frequency control, but it does not merge them into one large monolithic decision. Placement determines where the task should execute, using HEFT as a baseline and allowing controlled locality-driven deviations. DVFS then determines how aggressively the selected GPU can reduce frequency without violating criticality, backlog, and slack constraints.

\subsection{Optimization View}
\label{subsec:optimization_view}

At a high level, the runtime can be viewed as approximating the following constrained optimization objective:
\begin{align}
\min_{\pi,\Phi} \quad & J(E_{\mathrm{app}}, T_{\mathrm{app}}) \\
\mathrm{s.t.} \quad & T_{\mathrm{app}} \leq (1+\epsilon)T_{\mathrm{base}}, \\
& \pi(v) \in D, \quad \forall v \in V, \\
& \phi(v) \in \Phi, \quad \forall v \in V, \\
& \text{all DAG dependencies are respected},
\label{eq:runtime_optimization_problem}
\end{align}
where $\pi(v)$ is the placement decision, $\phi(v)$ is the frequency decision, and $J$ may be energy, EDP, ED$^2$P, or another user-selected objective. This formulation is a conceptual target rather than a claim that the online scheduler solves the constrained problem exactly.

Solving this problem exactly online is impractical. The thesis therefore uses a layered heuristic framework: HEFT-style baseline placement, controlled deviation for data locality, contextual-bandit adaptation of gate profiles, and proposal--commit DVFS for stable frequency control.

\subsection{Design Implications}
\label{subsec:design_implications}

The theory motivates four design choices:
\begin{enumerate}
    \item \textbf{Use the task graph.} Dependencies, data movement, and slack are essential for safe energy optimization.
    \item \textbf{Keep HEFT as a baseline.} A strong performance baseline reduces scheduling risk.
    \item \textbf{Adapt deviation online.} Workloads differ, so the scheduler should learn how aggressive it can be.
    \item \textbf{Commit DVFS by window.} Aggregating frequency intent avoids excessive transitions and unstable oscillations.
\end{enumerate}

\section{Summary}
\label{sec:theory_summary}

This chapter presented the theoretical foundation for energy-aware scheduling of multi-GPU task graphs. A CUDASTF application can be viewed as a DAG where tasks, dependencies, data residency, and communication costs are visible to the runtime. GPU DVFS introduces an energy--performance trade-off that depends on criticality, copy-wait behavior, backlog, frequency levels, and transition overhead. HEFT-style scheduling provides a strong baseline for performance, but it does not directly optimize energy or data-locality-driven trade-offs. Controlled deviation from HEFT, contextual bandit adaptation, region-level DVFS, and windowed proposal--commit frequency control form the theoretical basis for the framework developed in this thesis.
